%%%%%%%%%%%%Outline%%%%%%%%%%%%%%%%%%

%%%
% message<xx>
% xx

%%%%%%%%%%%%Outline%%%%%%%%%%%%%%%%%%

\section{Discussion}
\label{sec:discussion}

\textit{[TODO: Update discussion for v2 architecture. Key discussion
points to cover:]}

\textbf{LLVM as upper bound.} llvmbpf serves as an upper bound
reference, not a direct competitor. Running LLVM -O3 inside the kernel
is impractical; \textsc{BpfReJIT} captures a substantial portion of the
available headroom through a kernel-compatible framework.

\textbf{Complementarity with bytecode optimizers.} K2, Merlin, and
EPSO optimize BPF bytecode \emph{before} loading; \textsc{BpfReJIT}
optimizes programs \emph{after} loading, transparently and guided by
runtime data. A deployment could use both layers without interference,
as bytecode optimizations compose with post-load recompilation.

\textbf{Why userspace, not kernel.} Implementing all optimizations
in-kernel is infeasible for multiple reasons: the upstream review cycle
is measured in months to years; optimization algorithms require rapid
iteration that conflicts with kernel stability; workload-adaptive
decisions require runtime profiling data that the kernel JIT does not
collect; and fleet management capabilities (A/B testing, gradual
rollout, per-deployment customization) belong in userspace.

\textbf{Limitations.}
\textit{[TODO: Update limitations for v2. Include discussion of
prototype maturity, architecture coverage (x86 only so far), and
end-to-end deployment scope.]}